% TOP_PROBLEM: {"problemId": "problem.which-compact-convex-sets-are-realizable-as-limit-shapes-in-two-dimensional-first-passage-percolation", "canonicalTitle": "Realizable Limit Shapes in Two-Dimensional First-Passage Percolation", "releaseRank": 422, "primaryDomain": "probability_ergodic_dynamics", "primaryDomainLabel": "Probability, ergodic theory & dynamics", "displayStatus": "open", "releaseStatus": "open"}
% !TeX program = lualatex
% Definition and sources only; this file is not an LLM solution attempt.
\documentclass[11pt]{article}
\usepackage[margin=25mm]{geometry}
\usepackage{fontspec}
\setmainfont{Latin Modern Roman}
\usepackage{amsmath,amssymb,mathtools}
\usepackage{unicode-math}
\setmathfont{Latin Modern Math}
\usepackage{xurl,hyperref}
\hypersetup{colorlinks=true,urlcolor=blue}
\setcounter{secnumdepth}{0}
\setlength{\parindent}{0pt}
\setlength{\parskip}{5pt}
\setlength{\emergencystretch}{3em}
\begin{document}
\section*{\texorpdfstring{Realizable Limit Shapes in Two-Dimensional First-Passage Percolation}{Realizable Limit Shapes in Two-Dimensional First-Passage Percolation}}\label{422}
\subsection{Definitions and mathematical statement}
On nearest-neighbor edges of ${\mathbb{Z}}^2$, assign independent weights $t_e\ge0$ with law $\nu$. Passage time is $T(x,y)=\inf_{\gamma:x\to y}\sum_{e\in\gamma}t_e$. Assume $\nu(\{0\})<p_c(2)$ and ${\mathbb{E}}[\min(t_1,t_2,t_3,t_4)^2]<\infty$ for independent copies. The deterministic time-constant norm is obtained by the directional large-distance limit $\mu_\nu(x)=\lim_{n\to\infty}T(0,\lfloor nx\rfloor)/n$, extended continuously and homogeneously, and its unit ball is
\[
 B_\nu=\{x\in{\mathbb{R}}^2:\mu_\nu(x)\le1\}.
\]
These compact convex bodies describe the rescaled infected sets in the shape theorem. Characterize exactly which bodies occur as $\nu$ ranges over this class. Lattice reflection and coordinate-permutation symmetries are necessary, but are not asserted here to be sufficient. The minimum-moment condition must not be replaced by a stronger individual second-moment hypothesis.
\par\smallskip\textit{Formulation and concept references:} \href{https://sites.math.duke.edu/~rtd/FPP/50yrsFPP.pdf}{[S1]}.

\subsection{Short English statement}
Which deterministic planar shapes can arise from independent nearest-neighbor first-passage times satisfying the specified positivity and minimum-moment assumptions?
\subsection{Sources}
\begin{itemize}
\item[S1]Auffinger, Damron and Hanson, 50 years of first-passage percolation.
\url{https://sites.math.duke.edu/~rtd/FPP/50yrsFPP.pdf}.
\textit{Formulation source.}
\end{itemize}
\end{document}
